\subsection{Square Roots are Typically Irrational} 
\begin{theorem}
The square root of a positive integer is either a positive integer or an irrational number.
\end{theorem}
\begin{proof}
We'll argue that $\sqrt{2}$ is irrational. The same argument works for any prime number; a little more detail is needed for numbers that aren't prime.
\begin{itemize}
\item Suppose $\sqrt{2}$ were rational. Then it could be written as a fraction $\frac{p}{q}$ in lowest terms. That means that $p$ and $q$ have no \makeblank.
\item $\sqrt{2}=\dfrac{p}{q}$ means $2=\dfrac{p^2}{q^2}$. But $2=\dfrac{2}{1}=\dfrac{2q^2}{q^2}$, so $p^2=2q^2$.
\item That means \makeblanksmall is a factor of $p^2$. Since \makeblanksmall is prime, it must be a prime factor of $p$. Say $p=2r$.
\item Then $p^2=(2r)^2=2^2r^2$.
\item Going back to the fraction, we have $2=\dfrac{p^2}{q^2}=\dfrac{\makespace[0.5in]}{q^2}=\dfrac{\makespace[0.5in]}{q^2}=\dfrac{2q^2}{q^2}$.
\item Now we see that $2r^2=q^2$, so \makeblanksmall must be a factor of $q^2$ and thus $q$.
\item So 2 is a \makeblank[1.75in] of $p$ and $q$.
\item But at the beginning we picked $p$ and $q$ to have no \makeblank[1.75in]!
\item That means we can't pick values $p$ and $q$ like that, so $\sqrt{2}$ can't be written as a fraction---it is \makeblank.\qedhere
\end{itemize}
\end{proof}